\section{Installing the toolchain}\label{sec:00-setup:02-installing-toolchain:1}

Lean is managed by \textbf{elan}, a \textbf{version manager}. This is a small tool that installs and switches between different versions of Lean itself, so each project can pin the exact version it needs. (Readers familiar with \texttt{uv}'s Python-version management (\texttt{uv python install}) will recognize elan as playing the same role for Lean.)

\begin{enumerate}
\def\labelenumi{\arabic{enumi}.}
\item
  Install elan: follow the instructions at the official Lean installation guide, \href{https://lean-lang.org/lean4/doc/quickstart.html}{quickstart.html} (the ``Project setup'' page at \href{https://lean-lang.org/}{lean-lang.org} links the elan installer directly), or use a package manager. On Windows, Linux, and macOS alike, the recommended path is via VS Code's \emph{Lean 4} extension, which offers to install elan automatically through a platform-specific setup script for whichever OS it detects.
\item
  Verify installation:

\begin{lstlisting}[language=sh]
elan --version
lake --version
\end{lstlisting}
\item
  This repository ships a companion project, \texttt{../\allowbreak{}lean\_\allowbreak{}project}, already configured with:

\begin{lstlisting}
lean-toolchain: leanprover/lean4:v4.32.2
\end{lstlisting}

  Running \texttt{lake build} inside \texttt{lean\_\allowbreak{}project} causes elan to automatically download and use exactly that toolchain version.
\end{enumerate}
